\chapter{Problema 1: integrar um sistema já versionado ao GitLab institucional}
\label{cap:problema1}

Agora saímos da teoria. Este é um problema real, enfrentado no estágio.

\section{O cenário}

A instituição tem um GitLab empresarial restrito. Ao criar um repositório novo, ele já vem com:

\begin{itemize}
  \item uma branch \arq{main} criada;
  \item um \arq{README.md} padrão;
  \item regras institucionais;
  \item a \arq{main} \textbf{protegida}: não aceita \cmd{git push} direto, só Merge Request.
\end{itemize}

Só que vários sistemas \textbf{já existiam} nos computadores, e alguns já tinham anos de histórico Git próprio. Era preciso levar esses sistemas para o GitLab \textbf{sem perder o histórico}.

\section{Por que não funciona do jeito óbvio}

São dois históricos que nasceram separados:

\begin{center}
\begin{tikzpicture}
  \node[font=\small\bfseries\color{verdeMB}, anchor=west] at (-0.4,1.1) {REPOSITÓRIO LOCAL};
  \foreach \l/\x in {A/0,B/1.6,C/3.2,D/4.8}{\node[commit=verdeMB] (\l) at (\x,0.4) {\l};}
  \draw[hist=verdeMB] (A)--(B)--(C)--(D);
  \node[font=\small\bfseries\color{azulMB}, anchor=west] at (-0.4,-0.8) {REPOSITÓRIO GITLAB};
  \foreach \l/\x in {X/0,Y/1.6}{\node[commit] (\l) at (\x,-1.5) {\l};}
  \draw[hist] (X)--(Y);
  \node[font=\Large\color{vermelhoAlerta}] at (3.4,-1.5) {\faUnlink};
  \node[font=\small, anchor=west] at (4,-1.5) {nenhum ancestral em comum};
\end{tikzpicture}
\end{center}

O Git sabe juntar linhas que \textbf{nasceram de um mesmo ponto}. Entre \arq{A--B--C--D} e \arq{X--Y} não existe esse ponto. Por isso, um merge comum é recusado:

\begin{saida}
fatal: refusing to merge unrelated histories
\end{saida}

E um \cmd{git push} direto para a \arq{main} também é recusado, porque ela é protegida.

\section{O objetivo}

Transformar os dois históricos num só, sem perder nenhum commit:

\begin{center}
\begin{tikzpicture}
  \foreach \l/\x in {A/0,B/1.6,C/3.2,D/4.8}{\node[commit=verdeMB] (\l) at (\x,1) {\l};}
  \draw[hist=verdeMB] (A)--(B)--(C)--(D);
  \foreach \l/\x in {X/0,Y/2.4}{\node[commit] (\l) at (\x,-1) {\l};}
  \draw[hist] (X)--(Y);
  \node[commitm] (M) at (6.8,0) {M};
  \draw[hist=verdeMB] (D) to[out=0,in=135] (M);
  \draw[hist] (Y) to[out=0,in=-135] (M);
  \node[rotulo=verdeMB, anchor=west] at ([xshift=4mm]M.east) {desenvolvimento};
\end{tikzpicture}
\end{center}

O commit \textbf{M} é um commit de merge que une os dois históricos. Depois dele, a branch \arq{desenvolvimento} contém a \arq{main} institucional, então um Merge Request para a \arq{main} passa a funcionar normalmente.

\section{A solução, passo a passo}

\subsection*{Passo 1 — Adicionar o repositório remoto}

\begin{terminal}
git remote add origin <URL_DO_REPOSITORIO>
git remote -v
\end{terminal}
\begin{saida}
origin    <URL_DO_REPOSITORIO> (fetch)
origin    <URL_DO_REPOSITORIO> (push)
\end{saida}

Isso dá ao Git local um endereço remoto, chamado \arq{origin}. A URL aparece no botão \textbf{Code} (ou \textbf{Clone}) da página do projeto no GitLab.

\subsection*{Passo 2 — Buscar as informações do remoto}

\begin{terminal}
git fetch origin
\end{terminal}
\begin{saida}
From <URL_DO_REPOSITORIO>
 * [new branch]      main       -> origin/main
\end{saida}

O \cmd{fetch} baixa o histórico do GitLab para uma referência de leitura chamada \arq{origin/main}, \textbf{sem mexer} na sua branch nem nos seus arquivos.

\subsection*{Passo 3 — Renomear a branch local}

\begin{terminal}
git branch -m main desenvolvimento
\end{terminal}

A sua \arq{main} local vira \arq{desenvolvimento}. Isso evita confundi-la com a \arq{main} institucional, que é protegida e não pode receber push. Agora existem:

\begin{center}
\begin{tikzpicture}[l/.style={font=\small, anchor=west}]
  \node[rotulo=verdeMB] at (0,0) {desenvolvimento}; \node[l] at (1.8,0) {o histórico original do sistema (A--D)};
  \node[rotulo] at (-0.2,-0.8) {origin/main}; \node[l] at (1.8,-0.8) {o histórico institucional (X--Y)};
\end{tikzpicture}
\end{center}

\begin{dica}
Se a sua branch atual não se chama \arq{main}, use só \cmd{git branch -m desenvolvimento}, que renomeia a branch em que você está.
\end{dica}

\subsection*{Passo 4 — Unir os históricos}

\begin{terminal}
git merge --allow-unrelated-histories origin/main
\end{terminal}

O parâmetro \cmd{--allow-unrelated-histories} deve ser usado com consciência.

\begin{center}
\begin{tabularx}{\linewidth}{@{}XX@{}}
\begin{tcolorbox}[base, colback=vermelhoClaro, colframe=vermelhoAlerta, boxrule=0.8pt]
{\bfseries\color{vermelhoAlerta}\faTimes\; Ele NÃO significa}\\ \enquote{ignore todos os conflitos}.
\end{tcolorbox}
&
\begin{tcolorbox}[base, colback=verdeClaro, colframe=verdeMB, boxrule=0.8pt]
{\bfseries\color{verdeMB}\faCheck\; Ele significa}\\ \enquote{eu sei que esses históricos não têm ancestral comum e autorizo o Git a tentar integrá-los}.
\end{tcolorbox}
\end{tabularx}
\end{center}

Como os dois lados têm um \arq{README.md} diferente, surge um conflito:

\begin{saida}
Mesclagem automática de README.md
CONFLITO (adicionar/adicionar): conflito de mesclagem em README.md
Automatic merge failed; fix conflicts and then commit the result.
\end{saida}

\subsection*{Passo 5 — Resolver o conflito}

\begin{terminal}
git status
\end{terminal}
\begin{saida}
Caminhos não mesclados:
  (usar "git add <arquivo>..." para marcar resolução)
    ambos adicionaram:   README.md
\end{saida}

Aqui entra a parte que mais causa erro: \textbf{\texttt{ours} e \texttt{theirs}}.

\begin{center}
\begin{tikzpicture}
  \node[commitm, minimum size=13mm, font=\bfseries\color{azulMB}] (m) at (0,0) {merge};
  \node[caixa=verdeMB, text width=4.3cm] (o) at (-3.1,-2.2) {\textbf{OURS} (\enquote{nosso})\\ \small o lado da branch em que \textbf{você está} no momento do merge\\ \scriptsize aqui: \texttt{desenvolvimento}};
  \node[caixa, text width=4.3cm] (t) at (3.1,-2.2) {\textbf{THEIRS} (\enquote{deles})\\ \small o lado da branch que \textbf{está sendo incorporada}\\ \scriptsize aqui: \texttt{origin/main}};
  \draw[setafina] (m) -- (o); \draw[setafina] (m) -- (t);
\end{tikzpicture}
\end{center}

No nosso caso, queríamos manter o README do sistema, que está na branch em que estamos. Então:

\begin{terminal}
git restore --ours README.md
# forma tradicional, com o mesmo efeito:
git checkout --ours README.md
\end{terminal}

Se quiséssemos manter o README institucional:

\begin{terminal}
git restore --theirs README.md
# ou: git checkout --theirs README.md
\end{terminal}

\begin{atencao}[\texttt{ours} não é \enquote{o do meu computador}]
\texttt{ours} significa \textbf{o lado correspondente à branch atual durante aquele merge}. \texttt{theirs} é o outro lado. Os dois arquivos estão no seu computador. Interpretar errado faz você preservar a versão errada sem perceber.
\end{atencao}

\begin{atencao}[No rebase, os papéis se invertem]
Durante um \cmd{git rebase}, \texttt{ours} passa a ser a branch \emph{sobre a qual} você está reaplicando os commits, e \texttt{theirs} são os seus próprios commits. Nesta solução usamos \cmd{merge}, então vale a regra de cima. Se um dia usar rebase, confira antes.
\end{atencao}

\begin{dica}[E se eu quiser os dois READMEs?]
Você não precisa escolher um lado. Abra o \arq{README.md}, junte à mão as partes que interessam dos dois, apague os marcadores \texttt{<<<<<<<}, \texttt{=======} e \texttt{>>>>>>>}, e siga para o passo 6.
\end{dica}

\subsection*{Passo 6 — Registrar a resolução}

\begin{terminal}
git add README.md
git commit
\end{terminal}

Sem \cmd{-m}, o Git abre um editor com a mensagem \enquote{Merge remote-tracking branch 'origin/main' into desenvolvimento} já preenchida. Salve e feche. Este é o commit \textbf{M}: o histórico agora está integrado.

\begin{terminal}
git log --oneline --graph --all
\end{terminal}
\begin{saida}
*   89647ff Merge remote-tracking branch 'origin/main' into desenvolvimento
|\
| * 4738004 Initial commit
* 3ee51ff Ajuste
* 40d59e6 Primeira versão
\end{saida}

\subsection*{Passo 7 — Publicar a nova branch}

\begin{terminal}
git push origin desenvolvimento
\end{terminal}
\begin{saida}
 * [new branch]      desenvolvimento -> desenvolvimento
\end{saida}

No GitLab passam a existir, lado a lado, a \arq{main} (intocada e protegida) e a \arq{desenvolvimento} (o sistema com o histórico integrado).

\subsection*{Passo 8 — Abrir o Merge Request}

Este é o ponto final de todo o processo. Abra um Merge Request de \arq{desenvolvimento} para \arq{main}, pela interface web (capítulo \ref{cap:gitlab}) ou pelo terminal:

\begin{terminal}
glab mr create -s desenvolvimento -b main --title "Integra o sistema ao repositório institucional"
\end{terminal}

Como os dois históricos agora são relacionados e o conflito já foi tratado localmente, o merge \textbf{não é mais recusado}. A aprovação fica com quem tem permissão no repositório, conforme as regras do projeto.

\section{Resumo em oito linhas}

\begin{terminal}
git remote add origin <URL_DO_REPOSITORIO>          # 1. conectar
git fetch origin                                     # 2. baixar
git branch -m main desenvolvimento                   # 3. renomear
git merge --allow-unrelated-histories origin/main    # 4. unir
git restore --ours README.md                         # 5. decidir o conflito
git add README.md && git commit                      # 6. registrar
git push origin desenvolvimento                      # 7. publicar
glab mr create -s desenvolvimento -b main            # 8. pedir a integração
\end{terminal}

\section{Erros que podem aparecer}

\begin{center}
\renewcommand{\arraystretch}{1.4}
\small
\begin{tabularx}{\linewidth}{@{}>{\raggedright\arraybackslash\ttfamily\footnotesize}p{5.6cm} >{\raggedright\arraybackslash}X@{}}
\toprule
\textrm{\textbf{Mensagem}} & \textbf{O que fazer} \\
\midrule
fatal: refusing to merge unrelated histories & faltou o \cmd{--allow-unrelated-histories} no passo 4 \\
error: remote origin already exists. & o projeto já tem um \arq{origin}. Confira com \cmd{git remote -v}; para trocar a URL, use \cmd{git remote set-url origin <URL>} \\
remote: GitLab: You are not allowed to push code to protected branches on this project. & você tentou enviar para a \arq{main} protegida. Envie a \arq{desenvolvimento} e abra um Merge Request \\
error: Committing is not possible because you have unmerged files. & ainda há conflito sem resolver. Rode \cmd{git status}, resolva e faça \cmd{git add} \\
\bottomrule
\end{tabularx}
\end{center}

\begin{dica}
Quer treinar isso sem servidor nenhum? O exercício 4.1 (o desafio) do capítulo \ref{cap:exercicios} simula o GitLab institucional com uma pasta do seu computador e reproduz o caso inteiro.
\end{dica}
